% Source: https://artofproblemsolving.com/community/c3h1757873_list_of_logical_and_interesting_problems
% Stable identity: aops-topic-1757873-q04
The number has $9+2\cdot90+3=192$ digits, so $92$ remain. Compare numbers with the same number of digits from the left: we should make the first digit as large as possible, then the second, and so on, while keeping enough digits after each chosen one.

\begin{itemize}
\item The first $9$ appears after $8$ digits: delete $1,\dots,8$ ($8$ deletions).
\item The next $9$ is the last digit of $19$, after the $19$ digits $1,0,1,1,\dots,1,8,1$: delete them ($27$ deletions in all). Similarly take the $9$ of $29$, $39$ and $49$ ($46$, $65$, $84$ deletions).
\item Now $16$ deletions remain, and the digits continue $5,0,5,1,5,2,\dots,5,7,5,8,5,9,6,0,\dots$. The next $9$ (in $59$) would need $19$ deletions, too many. Among the next $17$ digits the largest is the $7$ of $57$, reached by deleting $15$ digits ($99$ in all).
\item With $1$ deletion left the digits continue $5,8,5,9,\dots$: delete the $5$ and keep the $8$ of $58$.
\end{itemize}
No deletions remain, so the rest, $5960616263\cdots99100$, is kept. The largest number is
\[
\boxed{9999978596061\cdots9899100},
\]
i.e.\ $9999978$ followed by the digits of $59,60,61,\dots,99,100$ (a $92$-digit number).
